\begin{lemma}
For every integer $t\geq2$ there is a constant $C=C(t)$ such that, for every
sufficiently large power of two $q$ and every $k\geq Cq(\operatorname{Nat.log}_2 q)^2$,
there is a loopless digraph $D$ (the paper's improved pair digraph
$D^*(t,q)$) that is $T_{t+1}$-free, has at least $q^{2t-1}/4$ vertices, and
satisfies
\[
 \operatorname{fwi}_k(D^*(t,q))\leq(Cq^t)^k.
\]
Here $\operatorname{Nat.log}_2$ is the base-two natural-number logarithm used
by the Lean threshold; for powers of two it is converted to the paper's
natural logarithm explicitly in the proof.
\end{lemma}
\begin{proof}
Fix $t\geq2$.  We choose one constant $C=C(t)$ sufficiently large below;
in particular, it is at least $16$ and all later lower bounds on $q$ are
absorbed into $q\geq C$.  Let $q\geq C$ be a power of two, write
$q=2^m$, and suppose
\[
        k\geq Cq(\operatorname{Nat.log}_2 q)^2.                 \tag{1}
\]
Since $q\geq2$, $\operatorname{Nat.log}_2 q=m$; put $L=m$.  Thus the exact
conversion between the two logarithms is
\[
        q=2^L,\qquad \log q=L\log2.
                                                               \tag{2}
\]
Consequently a bound $O_t(q\log q)$ in the paper is also $O_t(qL)$, with
the fixed factor $\log2$ absorbed into its constant.

Choose a field $K$ of order $q$ (for example, take the quotient of
$\mathbb F_2[X]$ by an irreducible polynomial of degree $m$).  Let
$G=G_K(t)$ be the projective polarity graph from
\noderef{PolarityGraph}: its vertices are the one-dimensional subspaces of
$K^{t+1}$ and $a\sim b$ means $a\perp b$.  The exact order, degree, and
spectral package supplied for this graph by
\noderef{OldPolarityParameters} is
\[
 n=|V(G)|=\frac{q^{t+1}-1}{q-1},\qquad
 d=\frac{q^t-1}{q-1},\qquad
 \lambda\leq2\sqrt d,
                                                               \tag{3}
\]
as well as
\[
 q^t/2\leq n\leq2q^t,\qquad
 q^{t-1}/2\leq d\leq2q^{t-1}.                                 \tag{4}
\]
The same node supplies the bilinear spectral inequality with these $n,d,
\lambda$ and the cardinality equality $|V(G)|=n$.  Applying
\noderef{ExpanderMixing} to that bilinear inequality and the indicator
vectors of arbitrary subsets gives the ordered estimate
\[
 \left|e_G(X,Y)-\frac dn|X||Y|\right|
 \leq\lambda\sqrt{|X||Y|}                                    \tag{5}
\]
for every pair of subsets $X,Y\subseteq V(G)$.  We will use (5) with its
ordered convention, including possible loops.

Define the witness digraph explicitly.  Its vertex set is
\[
 V(D^*)=\{(a,b)\in V(G)^2:a\perp b\},
\]
and, for two such ordered pairs, put
\[
 ((a,b),(a',b'))\in A(D^*)\quad\Longleftrightarrow\quad
 a\perp b'\ \text{ and }\ a'\not\perp b.                    \tag{6}
\]
This is the improved pair digraph described in the paper, rather than an
arbitrary existential digraph.  For each of the $n$ choices of $a$, the
number of choices of $b$ with $a\perp b$ is exactly $d$, including a
possible loop of $G$.  Therefore
\[
 |V(D^*)|=nd
 =\frac{q^{t+1}-1}{q-1}\frac{q^t-1}{q-1}.                    \tag{7}
\]
The geometric sums in (3) give $n\geq q^t$ and $d\geq q^{t-1}$, while
(4) gives the estimate needed uniformly below:
\[
 q^{2t-1}\leq |V(D^*)|\leq4q^{2t-1}.                         \tag{8}
\]
In particular, the required lower bound $q^{2t-1}/4$ holds, and the
quantity $4q^{2t-1}$, not $q^{2t-1}$, is a valid upper bound for the number
of children in the tuple tree.

The digraph is loopless directly from (6).  A loop at $(a,b)$ would require
both the vertex condition $a\perp b$ and the second arc condition
$a\not\perp b$, which is impossible.  Thus (6) defines a witness of the
loopless type used by \noderef{ForwardIndependentTuple} and
\noderef{ForwardIndependentCount}.

We next prove that $D^*$ is $T_{t+1}$-free.  Suppose that
$(a_i,b_i)_{i=1}^{t+1}$ were the image of a copy of the transitive
tournament in the sense of \noderef{TransitiveTournament}.  Choose nonzero
representatives $x_i,y_i\in K^{t+1}$ for $a_i,b_i$.  The vertex condition
and (6) give
\[
 \langle x_i,y_i\rangle=0,\qquad
 \langle x_i,y_j\rangle=0,
 \quad\langle x_j,y_i\rangle\neq0 \quad(i<j).                \tag{9}
\]
We claim that $y_1,\ldots,y_{t+1}$ are linearly independent.  If a nonzero
relation $\sum_i c_i y_i=0$ exists, let $i_0$ be its first index with
$c_{i_0}\neq0$.  If $i_0<t+1$, pair the relation with $x_{i_0+1}$.  All
terms with index below $i_0$ vanish by the choice of $i_0$; the diagonal
term at $i_0+1$ vanishes by (9); every term with index larger than
$i_0+1$ vanishes by the forward orthogonality in (9).  The only remaining
term is
\[
 c_{i_0}\langle x_{i_0+1},y_{i_0}\rangle,
\]
which is nonzero by the reverse nonorthogonality in (9), a contradiction.
If the first nonzero coefficient is instead at the final index $i_0=t+1$,
the relation is $c_{t+1}y_{t+1}=0$, impossible because
$c_{t+1}\neq0$ and the representative $y_{t+1}$ is nonzero.  This final-index
case is essential because there is no $x_{i_0+1}$ available.  Hence the
$y_i$ form a basis of $K^{t+1}$.

Now $x_1$ is orthogonal to $y_1$ by the vertex condition and to every
$y_j$ with $j>1$ by the forward arc condition in (9).  It is therefore
orthogonal to a basis.  The dot product on $K^{t+1}$ is nondegenerate, so
$x_1=0$, contradicting its choice as a representative.  This proves the
$T_{t+1}$-freeness asserted in the lemma.

It remains to count forward-independent tuples.  By the definitions in
\noderef{ForwardIndependentTuple} and
\noderef{ForwardIndependentCount}, a forward-independent tuple in $D^*$ is
equivalently a sequence
\[
 \sigma=(a_1,b_1,\ldots,a_m,b_m)
\]
with $a_i\perp b_i$ for every $i$ and, for every $i<j$,
\[
 a_i\perp b_j\ \Longrightarrow\ a_j\perp b_i.               \tag{10}
\]
Indeed, (10) is exactly the negation of the conjunction in the arc rule
(6).  This equivalence is reversible and also handles repeated entries,
because both sides use the same ordered conditions.

Form a rooted tree whose root is the empty sequence and whose children
append one pair satisfying (10).  Its level-$k$ vertices are counted by
$\operatorname{fwi}_k(D^*)$.  Fix a partial consistent sequence $\sigma$.
For each projective point $y$ define the linear subspace and its dimension
\[
 W^\sigma(y)=\mathop{\mathrm{span}}\{b_i:a_i\perp y\},\qquad
 r^\sigma(y)=\dim W^\sigma(y).                                \tag{11}
\]
If $(a,b)$ is a child of $\sigma$, then $a\perp b$.  Moreover, whenever
$a_i\perp b$, (10) gives $a\perp b_i$, so $a$ is orthogonal to all of
$W^\sigma(b)$.  The orthogonal complement has dimension $t+1-r(b)$;
therefore the number of possible first coordinates for a fixed $b$ is at
most
\[
 \frac{q^{t+1-r(b)}-1}{q-1}\leq2q^{t-r(b)}                 \tag{12}
\]
when $r(b)\leq t$, and is zero when $r(b)=t+1$.

For $0\leq r\leq t+1$, put
\[
 U_r^\sigma=\{y:r^\sigma(y)\leq r\},\qquad
 Z_r^\sigma=U_r^\sigma\setminus U_{r-1}^\sigma,\qquad
 U_{-1}^\sigma=\varnothing.                                  \tag{13}
\]
For each $0\leq r\leq t$ with $|U_r^\sigma|>0$, choose
$\ell=\ell_r^\sigma\in[0,r]$ for which $|Z_\ell^\sigma|$ is maximal.
Since the $Z_j^\sigma$ for $0\leq j\leq r$ partition $U_r^\sigma$,
\[
 |Z_\ell^\sigma|\geq\frac{|U_r^\sigma|}{r+1}
 \geq\frac{|U_r^\sigma|}{2t}.                                 \tag{14}
\]

Call $u\in Z_r^\sigma$ popular at level $r$ if it is contained in at
least $|Z_\ell^\sigma|/(16q)$ of the subspaces
$W^\sigma(y)$ with $y\in Z_\ell^\sigma$.  Each such subspace has dimension
$\ell$, hence contains
\[
 \frac{q^\ell-1}{q-1}\leq2q^{\ell-1}
\]
projective points (with the inequality trivial when $\ell=0).  Double
counting incidences $(u,y)$ with $u\subseteq W^\sigma(y)$, $u\in Z_r^\sigma$,
and $y\in Z_\ell^\sigma$, gives
\[
 \#\{\text{popular }u\}\,\frac{|Z_\ell^\sigma|}{16q}
 \leq |Z_\ell^\sigma|\,2q^{\ell-1},
 \qquad\text{so}\qquad
 \#\{\text{popular }u\}\leq32q^\ell.                         \tag{15}
\]
Mark every child whose second coordinate $b$ is popular for the level
$r=r(b)$.  By (12), (15), and $\ell_r^\sigma\leq r$, the number of marked
children of this type is at most
\[
 \sum_{r=0}^t 32q^{\ell_r^\sigma}\,2q^{t-r}
 \leq64(t+1)q^t\leq128tq^t.                                   \tag{16}
\]

Call a first coordinate $a$ poor at level $r$ if
\[
 |N_G(a)\cap Z_\ell^\sigma|
 \leq\frac{|Z_\ell^\sigma|}{8q},\qquad
 \ell=\ell_r^\sigma,                                          \tag{17}
\]
and let $P_r$ be the poor set.  The definition gives the first incidence
bound
\[
 e_G(P_r,Z_\ell^\sigma)\leq
 \frac{|P_r|\,|Z_\ell^\sigma|}{8q}.                             \tag{18}
\]
We also need the product bound in precisely the form of
\noderef{AlonRodlBound}.  From (4),
\[
 \frac dn\geq\frac1{4q},
\]
so (17) implies, for every $a\in P_r$,
\[
 |N_G(a)\cap Z_\ell^\sigma|\leq
 \frac d{2n}|Z_\ell^\sigma|.                                  \tag{19}
\]
The graph order and degree are positive, and (5) is exactly the all-subset
mixing premise required by \noderef{AlonRodlBound}.  Applying that node to
$A=P_r$ and $B=Z_\ell^\sigma$, with (19), yields
\[
 |P_r|\,|Z_\ell^\sigma|
 \leq\frac{4\lambda^2n^2}{d^2}.                               \tag{20}
\]
Using $\lambda^2\leq4d$, $n\leq2q^t$, and
$d\geq q^{t-1}/2$, the right side is at most $128q^{t+1}$; we retain the
weaker convenient estimate
\[
 |P_r|\,|Z_\ell^\sigma|\leq1024q^{t+1}.                        \tag{21}
\]
This is the Alon--Rodl product estimate, with its exact mixing hypotheses
made explicit rather than inferred from a proof of the cited node.

Since $|Z_r^\sigma|\leq|Z_\ell^\sigma|$, (5) and (21) give the second
incidence estimate
\[
\begin{aligned}
 e_G(P_r,Z_r^\sigma)
 &\leq\frac dn|P_r|\,|Z_r^\sigma|
       +\lambda\sqrt{|P_r|\,|Z_r^\sigma|}\\
 &\leq\frac4q\,1024q^{t+1}
       +4q^{(t-1)/2}\sqrt{1024q^{t+1}}\\
 &\leq5000q^t.                                                 \tag{22}
\end{aligned}
\]
Here $\lambda\leq2\sqrt d\leq4q^{(t-1)/2}$ follows from (3) and (4).
Mark every child with poor first coordinate $a$ at the level
$r=r(b)$.  A child also has $a\perp b$, and $r(b)=r$ puts
$b\in Z_r^\sigma$, so these children are counted by the ordered edges in
$e_G(P_r,Z_r^\sigma)$.  Summing (22) over $r$ marks at most
$5000(t+1)q^t\leq10000tq^t$ children of this second type.

Choose $A=A(t)$ large enough that it dominates the constants in (16) and
(22), and mark the root as well.  Every tree vertex then has at most
\[
                         h=Aq^t                              \tag{23}
\]
marked children.

We now prove the required bound on unmarked vertices on a path.  Suppose
that $(a,b)$ extends $\sigma$ and its child $\sigma'$ is unmarked.  Put
$r=r^\sigma(b)$ and $\ell=\ell_r^\sigma$.  The child exists only when
$r\leq t$, and $b\in Z_r^\sigma$.  Since $a$ is not poor, at least
$|Z_\ell^\sigma|/(8q)$ points $y\in Z_\ell^\sigma$ satisfy $a\perp y$.
Since $b$ is not popular, at most
$|Z_\ell^\sigma|/(16q)$ of these have $b\subseteq W^\sigma(y)$.
Consequently at least
\[
                         \frac{|Z_\ell^\sigma|}{16q}             \tag{24}
\]
points $y$ satisfy both $a\perp y$ and
$b\not\subseteq W^\sigma(y)$.  For each of them, (11) gives
\[
 W^{\sigma'}(y)=W^\sigma(y)+\langle b\rangle,\qquad
 r^{\sigma'}(y)=\ell+1.                                      \tag{25}
\]
Thus
\[
 |U_\ell^{\sigma'}|
 \leq |U_\ell^\sigma|-\frac{|Z_\ell^\sigma|}{16q}.             \tag{26}
\]
By (14), $|U_\ell^\sigma|\leq|U_r^\sigma|\leq2t|Z_\ell^\sigma|$.
Therefore every unmarked extension satisfies the quantitative decay
\[
 |U_\ell^{\sigma'}|\leq
 \left(1-\frac1{32tq}\right)|U_\ell^\sigma|.                  \tag{27}
\]
Record $\psi(\sigma')=\ell$.  Notice also that this record implies
$|U_\ell^\sigma|\geq|Z_\ell^\sigma|>0$, so the set being decreased was
nonempty immediately before every recorded decrease.

Consider a path
\[
 \sigma_0,\sigma_1,\ldots,\sigma_m
\]
from the root, and suppose it contains more than $w$ unmarked vertices
among $\sigma_1,\ldots,\sigma_m$.  Put
\[
                         w=AqL.                               \tag{28}
\]
By the pigeonhole principle, some $\ell\in[0,t]$ occurs as
$\psi(\sigma_i)$ for at least $w/(t+1)$ of those unmarked indices.  The
sets $U_\ell^{\sigma_i}$ are nonincreasing along the path, because adding
a pair can only enlarge each span in (11).  If $i^*$ is the last index
with this value of $\ell$, then (27), the preceding occurrences, and
$n\leq2q^t$ give
\[
 |U_\ell^{\sigma_{i^*-1}}|
 \leq n\left(1-\frac1{32tq}\right)^{w/(t+1)-1}.                \tag{29}
\]
Increase $A$ and $C$ if necessary.  Since $t+1\leq2t$, $q\geq C$, and
$w=AqL$, the elementary inequality $1-x\leq e^{-x}$ gives
\[
 \left(1-\frac1{32tq}\right)^{w/(t+1)-1}
 \leq \exp\left(-\frac{w}{100t^2q}\right)
 <\frac1{4q^t}.                                                \tag{30}
\]
The last inequality is ensured by taking $A>100t^2(t+1)$ and then taking
$C$ large enough that the harmless subtractive $1$ in the exponent is
absorbed.  Equations (29)--(30) imply
$|U_\ell^{\sigma_{i^*-1}}|<1$, hence this integer set is empty.  This
contradicts the nonemptiness recorded immediately before the unmarked
extension $\sigma_{i^*-1}\to\sigma_{i^*}$.  Thus every root-to-vertex path
has at most $w=AqL$ unmarked vertices.

Finally set
\[
 \Delta=4q^{2t-1},\qquad h=Aq^t,\qquad w=AqL.                  \tag{31}
\]
By (8), every tree vertex has at most $\Delta$ children.  By increasing
$C$ so that $q^{t-1}\geq A/4$, we have $h\leq\Delta$; by $L\geq1$ and
$C\geq A$, (1) gives $w\leq k$.  Let $\operatorname{count}(z)$ be the
number of level-$k$ paths whose Boolean signature records an unmarked
extension by true.  A signature with $j$ true entries has at most
$\Delta^j h^{k-j}$ paths, and (28)--(30) make this count zero for $j>w$.
The rooted-tree estimate of \noderef{RootedTreeCounting} therefore applies
with exactly the parameters in (31):
\[
 \operatorname{fwi}_k(D^*)
 \leq2^k\Delta^w h^{k-w}.                                     \tag{32}
\]
Because $q\geq4$, $4q^{2t-1}\leq q^{2t}$, and $h\geq1$, so
\[
 \operatorname{fwi}_k(D^*)
 \leq2^k q^{2tw}(Aq^t)^k.                                     \tag{33}
\]
The exact logarithm conversion (2) now gives
\[
 q^{2tw}=(2^L)^{2tAqL}=2^{2tAqL^2}\leq2^k,                   \tag{34}
\]
provided $C\geq2tA$, by (1).  Hence (33) is at most
$(4Aq^t)^k$, which is at most $(Cq^t)^k$ once $C\geq4A$.
This proves the forward-independent bound and completes the proof.
\end{proof}
